\section{Introduction}


The Russian invasion of Ukraine on 24 February 2022 initially generated strong bipartisan support in the United States. The Senate approved a \$40 billion aid package by 86 votes to 11 in May 2022, reflecting broad agreement about supporting a democratic ally under attack \citep{framWATCHSenateSends2022}. By late 2023 that agreement had dissolved. Ukraine aid featured prominently in Republican primary debates, and a government shutdown was avoided only after \$6 billion in assistance was stripped from it, days before the Speaker of the House was removed from office \citep{wangCandidatesDifferencesUkraine2023}. An issue that arrived without an established partisan polarization acquired one within eighteen months.

Party brand theory offers an account of why this might happen. Parties are understood to function as brands that signal positions to voters, and voters form stronger attachments to parties whose brands are distinguishable from one another \citep{lupuPartyBrandsPartisanship2013,lupuPartyBrandsCrisis2016}. Maintaining a brand requires two things at once: parties must occupy distinct positions, and each must maintain internal consistency in positions. An issue without a settled partisan meaning therefore presents an opportunity, particularly for the party that does not already own it. The transformation of Ukraine aid from consensus to contestation is a naturally occurring case in which both conditions can be observed forming rather than inferred from an established polarization.

Observing them requires measuring how parties were represented rather than what they said. Brands are shaped by indirect communication as much as by direct communication, since most voters encounter party positions through journalists who select which statements to report and how to frame the conflict between them \citep{karlbergPartisanBrandingMedia2002,levenduskyDoesMediaCoverage2016}. Existing work on party brands has drawn largely on roll-call votes and floor speech, which capture official positions but not how those positions reached publics. We therefore take news coverage as the primary source of measurement.

We retrieved United States news articles published between March 2022 and December 2024, extracted sentences that discuss Ukraine aid and are attributable to one party, and annotated each on four ideological dimensions using a large language model. Party and politician references are masked before annotation, so that positions are recovered from what a sentence conveys rather than from prior association with a party. We then aggregated sentence-level annotations to articles and articles to months, producing continuous series of between-party polarization and within-party message coherence across 34 months.

We organize the investigation around four questions. How did polarization on general Ukraine aid stance develop across the period, and did it accumulate gradually or shift in stages? Which party contributed more to the widening polarization, and which of three explanatory dimensions moved most closely with it? How is polarization related to each party's within-party message coherence? And how did a combined index of brand conditions evolve?

This study makes three types of contribution. Theoretically, we separate two conditions that party brand theory treats together. The theory holds that a usable brand requires both distinctiveness and consistency, and empirical work has generally assumed that a party consolidating a distinctive position will also come to speak about it more consistently. We show that the two can advance at different rates within the same party: Republican coverage moved further from the Democratic position and became more internally consistent across the period, yet the two did not move together from month to month. We also extend the asymmetric incentives literature from positional movement to internal consolidation, since the party that moved was also the only party whose messaging tightened, while the party defending an established commitment did neither. And we show that brand formation here proceeded in stages rather than gradually, with a period of elevated polarization that did not hold followed by one that did.

Empirically, we provide the first continuous measurement of how the two parties were represented on Ukraine aid across the full period from March 2022 to December 2024. Prior work on this case has measured mass discourse at a single early moment, elite speech through qualitative interpretation, or policy rationales without systematic measurement. Our series covers 37,312 annotated sentences drawn from 11,819 articles across 34 months, which allows the timing and structure of change to be recovered. We also report two findings that aligns with previous literature: the framing that set foreign commitments against domestic priorities aligned most closely with the aid stance polarization, and neither party's coherence contemporaneously tracked the corresponding polarization.

Methodologically, we introduce the Party Brand Index, a monthly measure combining between-party divergence with each party's within-party coherence through a geometric mean. Because a geometric mean is dampened whenever its components develop unevenly, the index detects joint development rather than summarizing divergence, which makes answerable a question the theory poses and measures of divergence alone cannot address. We also set out an annotation design for recovering positions on multiple predetermined dimensions from single sentences, validate it on three properties, and treat measurement uncertainty jointly by resampling whole articles once and recomputing every quantity from each replicate.

The remainder of the paper proceeds as follows. We review literature on party brands, foreign policy polarization, and computational approaches to measuring both, and set out our research questions and hypotheses. We then describe the corpus, the annotation and its validation, the construction of our measures, and the temporal models. Results are reported by research question. We close by discussing what the findings imply for party brand theory and for debates about the durability of foreign policy consensus.

